% =====================================================================
\chapter{Variables aleatorias discretas}
\label{tema:2}
% =====================================================================

Una variable aleatoria asigna un número a cada resultado de un experimento. Este tema estudia las variables discretas, entre ellas la distribución categórica que produce un LLM en cada paso, y sus dos resúmenes principales, la esperanza y la varianza. Se sigue el capítulo~2 y la §4.2 de~\cite{bt}, las lecciones~5 a~7 de~\cite{mit6041} y las clases~4, 5a y~7 de~\cite{mit1805}; la orientación al aprendizaje automático procede de~\cite[§6.2 y §6.4]{mml}.

\section{Variables aleatorias y funciones de masa}

\begin{definicion}[Variable aleatoria discreta]\label{def:va}
Una \emph{variable aleatoria} es una función $X\colon\Omega\to\R$ tal que $\{\omega: X(\omega)\le x\}\in\mathcal{F}$ para todo $x\in\R$. Es \emph{discreta} si toma un conjunto finito o numerable de valores. Su \emph{función de masa} es
\[
p_X(x)=\PP(X=x),
\]
que cumple $p_X(x)\ge 0$ y $\sum_x p_X(x)=1$. Si $X$ y $Y$ tienen la misma función de masa se dice que tienen la misma \emph{distribución}.
\end{definicion}

\begin{definicion}[Distribuciones discretas básicas]\label{def:discretas}
\begin{enumerate}
  \item $X\sim\Ber(p)$, con $p\in[0,1]$, si $X\in\{0,1\}$ y $p_X(x)=p^{x}(1-p)^{1-x}$. Modela un experimento con dos resultados: acierto o fallo.
  \item $X\sim\Bin(n,p)$ si $X$ es el número de éxitos en $n$ ensayos de Bernoulli independientes con probabilidad $p$:
  \[
  p_X(k)=\binom{n}{k}p^{k}(1-p)^{n-k},\qquad k=0,1,\dots,n.
  \]
  \item $X\sim\Cat(p_1,\dots,p_K)$, con $p_k\ge 0$ y $\sum_k p_k=1$, si $X\in\{1,\dots,K\}$ y $\PP(X=k)=p_k$. Equivalentemente, $p_X(k)=\prod_{j=1}^{K}p_j^{\mathbf{1}[j=k]}$.
\end{enumerate}
\end{definicion}

La categórica es la generalización de la Bernoulli a $K$ resultados, y es exactamente lo que produce un LLM en cada posición: la softmax convierte los $K$ logits en el vector $(p_1,\dots,p_K)$ de una categórica sobre el vocabulario. Los valores $1,\dots,K$ son etiquetas sin orden ni distancia: el token~517 no está «entre» el~516 y el~518.

\begin{ejemplo}
Que la función de masa binomial sume~1 es el desarrollo del binomio de Newton: $\sum_{k=0}^{n}\binom nk p^{k}(1-p)^{n-k}=(p+(1-p))^{n}=1$. El coeficiente $\binom nk$ cuenta las secuencias de $n$ ensayos con exactamente $k$ éxitos, y cada una tiene probabilidad $p^{k}(1-p)^{n-k}$ por la independencia de los ensayos~\cite[§2.2]{bt}.
\end{ejemplo}

\section{Esperanza}

\begin{definicion}[Esperanza]\label{def:esperanza}
La \emph{esperanza} o \emph{media} de una variable discreta $X$ es
\[
\E[X]=\sum_x x\,p_X(x),
\]
siempre que la serie converja absolutamente.
\end{definicion}

La esperanza es el centro de gravedad de la función de masa. Para una categórica con etiquetas arbitrarias no tiene sentido «el token medio», pero sí la esperanza de cualquier función numérica de los tokens, por ejemplo de la pérdida.

\begin{teorema}[Esperanza de una función]\label{teo:lotus}
Si $g\colon\R\to\R$ y $X$ es discreta, entonces $\E[g(X)]=\sum_x g(x)\,p_X(x)$.
\end{teorema}

\begin{proof}
Sea $Y=g(X)$. Para cada valor $y$, $p_Y(y)=\sum_{x:g(x)=y}p_X(x)$. Entonces $\E[Y]=\sum_y y\,p_Y(y)=\sum_y\sum_{x:g(x)=y}g(x)\,p_X(x)=\sum_x g(x)\,p_X(x)$, porque los conjuntos $\{x:g(x)=y\}$ forman una partición de los valores de $X$.
\end{proof}

\begin{proposicion}[Linealidad]\label{prop:linealidad}
Para variables $X$, $Y$ con esperanza finita y constantes $a,b,c$,
\[
\E[aX+bY+c]=a\,\E[X]+b\,\E[Y]+c.
\]
La igualdad no requiere independencia.
\end{proposicion}

\begin{proof}
Sea $p_{X,Y}$ la función de masa conjunta (\cref{def:conjunta}). Por la versión bidimensional del \cref{teo:lotus},
$\E[aX+bY+c]=\sum_{x,y}(ax+by+c)\,p_{X,Y}(x,y)=a\sum_x x\sum_y p_{X,Y}(x,y)+b\sum_y y\sum_x p_{X,Y}(x,y)+c$, y las sumas interiores son las funciones de masa marginales.
\end{proof}

\begin{ejemplo}[Dado equilibrado]\label{ej:dado}
Si $X$ es la puntuación de un dado equilibrado, $\E[X]=(1+2+\dots+6)/6=21/6=3{,}5$. Por el \cref{teo:lotus}, $\E[X^{2}]=(1+4+9+16+25+36)/6=91/6$.
\end{ejemplo}

\section{Varianza}

\begin{definicion}[Varianza]\label{def:varianza}
La \emph{varianza} de $X$ es $\Var(X)=\E\bigl[(X-\E[X])^{2}\bigr]$, y su \emph{desviación típica} es $\sigma_X=\sqrt{\Var(X)}$, que tiene las mismas unidades que $X$.
\end{definicion}

\begin{proposicion}\label{prop:varianza}
\begin{enumerate}
  \item $\Var(X)=\E[X^{2}]-\E[X]^{2}$.
  \item $\Var(aX+b)=a^{2}\Var(X)$ para constantes $a,b$.
  \item $\Var(X)\ge 0$, con igualdad si y solo si $X$ es constante con probabilidad~1.
\end{enumerate}
\end{proposicion}

\begin{proof}
Sea $\mu=\E[X]$. (i) Desarrollando el cuadrado y por linealidad, $\E[(X-\mu)^{2}]=\E[X^{2}]-2\mu\E[X]+\mu^{2}=\E[X^{2}]-\mu^{2}$. (ii) $\E[aX+b]=a\mu+b$, luego $\Var(aX+b)=\E[(aX-a\mu)^{2}]=a^{2}\E[(X-\mu)^{2}]$. (iii) $\Var(X)=\sum_x(x-\mu)^{2}p_X(x)$ es una suma de términos no negativos; es nula si y solo si $p_X(x)=0$ para todo $x\ne\mu$.
\end{proof}

\begin{ejemplo}\label{ej:tresvalores}
Para el dado del \cref{ej:dado}, $\Var(X)=91/6-(7/2)^{2}=182/12-147/12=35/12\approx 2{,}917$. Si $X\in\{0,1,2\}$ con probabilidades $(0{,}5;\,0{,}3;\,0{,}2)$, entonces $\E[X]=0{,}3+0{,}4=0{,}7$, $\E[X^{2}]=0{,}3+0{,}8=1{,}1$ y $\Var(X)=1{,}1-0{,}49=0{,}61$.
\end{ejemplo}

Para $X\sim\Ber(p)$, $X^{2}=X$, luego $\E[X]=\E[X^{2}]=p$ y $\Var(X)=p-p^{2}=p(1-p)$, máxima en $p=1/2$.

\section{Varias variables aleatorias}

\begin{definicion}[Distribución conjunta e independencia]\label{def:conjunta}
La \emph{función de masa conjunta} de $X$ e $Y$ es $p_{X,Y}(x,y)=\PP(X=x,\,Y=y)$; las \emph{marginales} se obtienen sumando, $p_X(x)=\sum_y p_{X,Y}(x,y)$. Las variables $X$ e $Y$ son \emph{independientes} si $p_{X,Y}(x,y)=p_X(x)\,p_Y(y)$ para todo $(x,y)$, y $X_1,\dots,X_n$ son independientes si la conjunta factoriza en el producto de las $n$ marginales.
\end{definicion}

\begin{definicion}[Covarianza]
La \emph{covarianza} de $X$ e $Y$ es $\Cov(X,Y)=\E\bigl[(X-\E[X])(Y-\E[Y])\bigr]=\E[XY]-\E[X]\,\E[Y]$. Si $\Cov(X,Y)=0$, las variables se llaman \emph{incorreladas}.
\end{definicion}

\begin{proposicion}\label{prop:suma}
\begin{enumerate}
  \item Si $X$ e $Y$ son independientes, $\E[XY]=\E[X]\,\E[Y]$ y, por tanto, $\Cov(X,Y)=0$.
  \item $\Var(X+Y)=\Var(X)+\Var(Y)+2\Cov(X,Y)$. En particular, la varianza de una suma de variables independientes es la suma de las varianzas.
\end{enumerate}
\end{proposicion}

\begin{proof}
(i) Por la factorización de la conjunta, $\E[XY]=\sum_{x,y}xy\,p_X(x)p_Y(y)=\bigl(\sum_x x\,p_X(x)\bigr)\bigl(\sum_y y\,p_Y(y)\bigr)$. (ii) Sean $\tilde X=X-\E[X]$ e $\tilde Y=Y-\E[Y]$. Entonces $\Var(X+Y)=\E[(\tilde X+\tilde Y)^{2}]=\E[\tilde X^{2}]+\E[\tilde Y^{2}]+2\E[\tilde X\tilde Y]$.
\end{proof}

El recíproco de (i) es falso: hay variables incorreladas que no son independientes (\cref{eje:2-covarianza}).

\begin{corolario}[Varianza de la media muestral]\label{cor:media}
Sean $X_1,\dots,X_n$ independientes con media $\mu$ y varianza $\sigma^{2}$, y sea $\Xbar_n=\frac1n\sum_{i=1}^{n}X_i$. Entonces $\E[\Xbar_n]=\mu$ y $\Var(\Xbar_n)=\sigma^{2}/n$.
\end{corolario}

\begin{proof}
La media se obtiene por linealidad. Para la varianza, por la \cref{prop:varianza}(ii) y la \cref{prop:suma}(ii) aplicada repetidamente, $\Var(\Xbar_n)=\frac{1}{n^{2}}\sum_{i}\Var(X_i)=\frac{n\sigma^{2}}{n^{2}}$.
\end{proof}

El \cref{cor:media} es el resultado central del módulo: promediar $n$ observaciones independientes divide la varianza por $n$ y la desviación típica por $\sqrt n$. El tema~5 lo convierte en garantías sobre el error de cualquier promedio. Como aplicación inmediata, si $X\sim\Bin(n,p)$ se escribe $X=\sum_{i=1}^{n}B_i$ con $B_i\sim\Ber(p)$ independientes, y se obtiene $\E[X]=np$ y $\Var(X)=np(1-p)$.

\begin{table}[h]
\centering\small
\caption{Distribuciones de los temas~2 y~3 y uso en un LLM.}\label{tab:distribuciones}
\begin{tabular}{@{}lllll@{}}
\toprule
Distribución & Valores & $\E[X]$ & $\Var(X)$ & Uso en un LLM \\
\midrule
$\Ber(p)$ & $\{0,1\}$ & $p$ & $p(1-p)$ & acierto en una pregunta de evaluación \\
$\Bin(n,p)$ & $\{0,\dots,n\}$ & $np$ & $np(1-p)$ & aciertos en un banco de pruebas \\
$\Cat(p_1,\dots,p_K)$ & $\{1,\dots,K\}$ & --- & --- & token siguiente \\
$\Unif(a,b)$ & $[a,b]$ & $(a+b)/2$ & $(b-a)^{2}/12$ & número aleatorio del muestreo \\
$\Expo(\lambda)$ & $[0,\infty)$ & $1/\lambda$ & $1/\lambda^{2}$ & tiempos de espera \\
$\N(\mu,\sigma^{2})$ & $\R$ & $\mu$ & $\sigma^{2}$ & pesos iniciales, ruido \\
\bottomrule
\end{tabular}
\end{table}

\section{Muestreo de una distribución categórica}

Para generar un token, el LLM necesita extraer una muestra de $\Cat(p_1,\dots,p_K)$ a partir de un número aleatorio uniforme. Sea $F_k=p_1+\dots+p_k$ la probabilidad acumulada, con $F_0=0$ y $F_K=1$.

\begin{proposicion}[Transformada inversa discreta]\label{prop:inversa-discreta}
Sea $U\sim\Unif(0,1)$, es decir, $\PP(a<U\le b)=b-a$ para $0\le a\le b\le 1$ (\cref{tema:3}). La variable
\[
X=\min\{k: U\le F_k\}
\]
tiene distribución $\Cat(p_1,\dots,p_K)$.
\end{proposicion}

\begin{proof}
Como $F_0\le F_1\le\dots\le F_K$, se tiene $X=k$ si y solo si $F_{k-1}<U\le F_k$. Por tanto, $\PP(X=k)=F_k-F_{k-1}=p_k$.
\end{proof}

\begin{figure}[h]
\centering
\begin{tikzpicture}[x=10cm,y=1cm,font=\small]
  \draw[thick] (0,0) -- (1,0);
  \foreach \x/\l in {0/0,0.5/{0{,}5},0.8/{0{,}8},1/1} {\draw (\x,0.12) -- (\x,-0.12) node[below] {$\l$};}
  \draw[-{Stealth[length=2mm]},thick] (0.6,0.9) node[above] {$u=0{,}6$} -- (0.6,0.08);
  \draw[decorate,decoration={brace,amplitude=4pt,mirror},yshift=-0.55cm] (0,0) -- (0.5,0) node[midway,below=4pt,align=center] {token 0\\ $p_0=0{,}5$};
  \draw[decorate,decoration={brace,amplitude=4pt,mirror},yshift=-0.55cm] (0.5,0) -- (0.8,0) node[midway,below=4pt,align=center] {token 1\\ $p_1=0{,}3$};
  \draw[decorate,decoration={brace,amplitude=4pt,mirror},yshift=-0.55cm] (0.8,0) -- (1,0) node[midway,below=4pt,align=center] {token 2\\ $p_2=0{,}2$};
\end{tikzpicture}
\caption{Muestreo por transformada inversa del \cref{ej:muestreo}: el intervalo $[0,1]$ se divide en tramos de longitud $p_k$ y se devuelve el tramo en el que cae $u$.}
\label{fig:inversa}
\end{figure}

\begin{ejemplo}[Muestreo de un token]\label{ej:muestreo}
Un vocabulario de tres tokens, numerados $0$, $1$ y $2$, recibe las probabilidades $(0{,}5;\,0{,}3;\,0{,}2)$. Las acumuladas son $(0{,}5;\,0{,}8;\,1)$. Si el generador produce $u=0{,}6$, como $0{,}5<0{,}6\le 0{,}8$, se emite el token~1 (\cref{fig:inversa}).
\end{ejemplo}

\begin{observacion}[Conexión con los LLM]
La softmax con temperatura del módulo~1 modifica el vector $(p_k)$ antes de aplicar la \cref{prop:inversa-discreta}: con $T\to 0$ toda la masa se concentra en el logit máximo (si es único) y el muestreo se vuelve determinista; con $T$ grande, la categórica se acerca a la uniforme. Por eso la misma instrucción puede producir respuestas distintas en dos ejecuciones: cada respuesta es una realización distinta de la misma secuencia de distribuciones categóricas.
\end{observacion}

\begin{lecturas}
\item Bertsekas y Tsitsiklis~\cite{bt}, §2.1–2.4 (variables discretas, funciones de masa, esperanza y varianza), §2.5 y §2.7 (conjuntas e independencia) y §4.2 «Covariance and Correlation».
\item MIT 6.041SC~\cite{mit6041}, unidad~I, lecciones~5 «Discrete Random Variables; Probability Mass Functions; Expectations», 6 y~7.
\item MIT 18.05~\cite{mit1805}, clases~4a–4b (variables discretas y valor esperado), 5a (varianza) y 7b «Covariance and Correlation».
\item Deisenroth, Faisal y Ong~\cite{mml}, §6.2 «Discrete and Continuous Probabilities» y §6.4 «Summary Statistics and Independence».
\end{lecturas}

\section*{Ejercicios del tema 2}
\addcontentsline{toc}{section}{Ejercicios del tema 2}

\begin{ejercicio}{1}{eje:2-categorica}
¿Por qué la salida de un LLM se modela con una distribución categórica y no con una normal sobre el índice del token? Dé al menos dos razones.
\end{ejercicio}

\begin{ejercicio}{2}{eje:2-dados}
Se lanzan dos dados equilibrados independientes. Calcule la esperanza y la varianza de la suma $S$ y la probabilidad $\PP(S=7)$.
\end{ejercicio}

\begin{ejercicio}{2}{eje:2-softmax}
Un modelo con cuatro tokens, numerados de~0 a~3, produce los logits $z=(2,1,0,-1)$. Calcule la softmax con tres decimales, las probabilidades acumuladas y el token emitido según la \cref{prop:inversa-discreta} cuando $u=0{,}5$ y cuando $u=0{,}9$.
\end{ejercicio}

\begin{ejercicio}{2}{eje:2-dados4}
Se lanzan dos dados de cuatro caras (numeradas del~1 al~4). Cada uno de los 16 resultados ordenados $(i,j)$ tiene probabilidad proporcional a la suma $i+j$. (a) Halle la constante de proporcionalidad. (b) Calcule la probabilidad de que la suma sea par. (c) Calcule la probabilidad de que un dado muestre un~2 y el otro un~3. (d) Calcule la esperanza de la suma.
\fuente{adaptado de MIT 6.041~\cite{mit6041f10}, \emph{Problem Set}~1, problema~3 (el apartado (d) es nuevo).}
\end{ejercicio}

\begin{ejercicio}{2}{eje:2-binomial}
Un modelo responde a 10 preguntas independientes y acierta cada una con probabilidad $0{,}8$. Sea $X$ el número de aciertos. Calcule $\E[X]$, $\Var(X)$, $\PP(X=10)$ y $\PP(X\ge 9)$.
\end{ejercicio}

\begin{ejercicio}{3}{eje:2-covarianza}
(a) Demuestre la \cref{prop:suma} con todo detalle, justificando cada uso de la linealidad. (b) Sea $X$ uniforme en $\{-1,0,1\}$ e $Y=X^{2}$. Demuestre que $X$ e $Y$ son incorreladas pero no independientes.
\end{ejercicio}
